\section{Introduction}

In recent years, large language models (LLMs) \citep{achiam2023gpt, chowdhery2023palm, touvron2023llama} have achieved impressive breakthroughs in natural language understanding and generation, marking a critical advancement in the exploration of artificial general intelligence (AGI). 
As the capabilities of LLMs progress, LLM-empowered agents~\citep{qin2023toolllm, dong2024self, wang2024survey} emerge as key components for integrating expertise and tools to effectively translate these advancements into practical applications.
Advancing this paradigm, multi-agent systems~\citep{shen2024hugginggpt, agentscope, wu2023autogen}, which involve multiple diverse agents, provide great flexibility and adaptability by leveraging and collaborating on the strengths of various agents, promoting comprehensive solutions to complex real-world problems.

Previous studies~\citep{cai2023low, qian2023communicative} propose to define suitable standard operating procedures (SOPs) based on the insights and experiences of human professionals for solving specific tasks, such as software development \citep{hong2023metagpt} and simulating personality traits \citep{serapio2023personality}. In these scenarios, multiple agents are assigned to execute tasks following the predefined SOPs, resulting in remarkable successes. However, for multi-agent systems designed to tackle a diverse range of complex real-world problems, there is a critical need for a central entity that can automatically generate task-specific operating procedures and coordinate the activities of various agents \citep{wang2024survey}.

This central entity, typically referred to as a meta-agent (a.k.a. a controller or planner), carries two primary responsibilities. Firstly, the meta-agent needs to comprehend user queries and decompose them into several sub-tasks, ensuring that each sub-task can be adequately addressed by a single agent. Secondly, the meta-agent is also expected to assign these sub-tasks to the appropriate agents for execution, enabling the solutions of these sub-tasks to collectively provide comprehensive answers to the original user queries.


However, since the meta-agent cannot effectively associate sub-tasks with agents based on the provided agents descriptions \citep{liu2024dynamic, hong2023metagpt}, the performance of task decomposition and allocations might be sub-optional. To tackle these challenges, we identify three design principles for an agent-oriented planning framework, including solvability, completeness, and non-redundancy. These principles further inspire us to propose {\bf \name}, a novel framework for {\bf A}gent-{\bf O}riented {\bf P}lanning in multi-agent systems.


Specifically, to resolve a user query, the meta-agent first performs fast task decomposition and allocation, serving as intermediate results that can be further modified and revised. The proposed \name incorporates a reward model designed to efficiently evaluate the solvability of sub-tasks without requiring actual agent calls.
According to the evaluation results, some sub-tasks may be executed by the assigned agents, others may be deemed inappropriate and require re-planning, while the remaining sub-tasks are further assessed using the representative works mechanism, which helps determine whether they should be planned in detail or re-described for better aligning the ability of agents.
Meanwhile, a detector is utilized to identify the missing key information or redundant content in the decomposed sub-tasks, and to provide suggestions to the meta-agent for enhancing the completeness and non-redundancy.
Besides, a feedback loop is integrated into \name to promote ongoing enhancements for the problem-solving process.

Extensive experiments are conducted based on several reasoning datasets that require collaboration among multiple LLM-empowered agents. Comparisons between \name and baseline methods demonstrate the remarkable advancements achieved by the proposed framework. 
Furthermore, we conduct an ablation study to show the contributions of different components in \name, and provide discussions on the potential for further enhancing agent-oriented planning within multi-agent systems.



